\section{Related literature}\label{sec:literature}

The paper connects four literatures.

\paragraph{The economics of stablecoins}
A first strand studies what keeps stablecoins at par and what they do to the rest of the financial system.
\citet{lyons2023stable} show how arbitrage capacity and the issuer's balance sheet determine deviations from the peg;
\citet{gorton2023wildcat} draw the analogy with free-banking-era private money; \citet{arner2020stablecoins} and the
\citet{bis2025aer} discuss risks and regulation, and \citet{liao2022stablecoins} the implications for bank funding.
Work on the interaction between stablecoin issuance and crypto prices includes \citet{griffin2020tether},
\citet{ante2021issuance} and \citet{kristoufek2021tethered}. This literature treats stablecoins mostly as a single
asset. I show that the same token on different blockchains serves different clienteles---USDT on Tron is a small-balance,
high-turnover rail, USDT on Ethereum a large-balance one---which matters for any analysis of who uses stablecoins and why.

\paragraph{Illicit finance in crypto-assets}
\citet{foley2019sex} estimate the share of bitcoin activity associated with illegal goods using network clustering.
\citet{cong2023anatomy} provide an anatomy of crypto-enabled cybercrime, and \citet{griffin2024slavery} trace the flows
of ``pig-butchering'' fraud from victims to exchanges, documenting the role of USDT and large centralised exchanges in
cashing out. Policy reports by \citet{unodc2024casinos,unodc2025inflection} and the
\citet{fatf2020redflags,fatf2023cyberfraud,fatf2025targeted} describe typologies and red flags. Industry reports such as
\citet{chainalysis2025crime} provide aggregate estimates whose attribution methods are proprietary. Relative to this
work I focus on a single jurisdiction at the boundary between the banking system and on-chain rails, and on the
identification problem: how to distinguish laundering from the legitimate motives that produce similar data.

\paragraph{Market structure and data quality}
\citet{makarov2020arbitrage} document persistent cross-country price premia in crypto markets, such as the Korean
``kimchi premium'', that reflect capital controls; I use the analogous USDT/THB premium as a fingerprint of the
store-of-value motive. \citet{cong2023wash} show that reported exchange volumes are inflated by wash trading, which is
why the plan relies on ratios within a platform and on adjusted volumes rather than on raw cross-exchange comparisons.
\citet{auer2022retail} relate retail adoption to prices using app-download data.

\paragraph{Blockchain forensics and anomaly detection}
Address clustering and labelling go back to \citet{meiklejohn2013fistful}; supervised and graph-based approaches to
anti-money-laundering on transaction graphs include \citet{weber2019aml}. I use unsupervised scores---the isolation forest
of \citet{liu2008iforest} and a transparent composite of red flags---because labels are scarce in practice, and I show
formally and in simulation that the direction of fund flows must be respected when exposure to flagged addresses is
measured. Community detection, used in the replication package for exploratory clustering, follows
\citet{blondel2008louvain}.

\paragraph{Methods}
The bunching estimator adapts the notch and kink designs of \citet{saez2010bunch}, \citet{chetty2011adjustment} and
\citet{kleven2013notches}, surveyed by \citet{kleven2016bunching}, to a setting in which the notch is a verification
step rather than a tax, and in which a within-platform comparison across assets removes the honest response. Policy
evaluation relies on synthetic control \citep{abadie2010synth,abadie2021using} and, as a robustness check, synthetic
difference-in-differences \citep{arkhangelsky2021sdid}. Cross-country controls include the capital-account openness
index of \citet{chinn2006capital}.
